% =====================================================================
%  Chapter 10  ·  Session 10 --- Logging and log integrity
% =====================================================================
\chapter{Logging and log integrity}

Session~9 observed the signals; this chapter asks whether the record of
them can be believed. Logs have two roles. Operationally they show how a
system behaves; in security they are evidence, namely the timeline of an
attack, the record of who did what, the material an incident responder, an
auditor or a court relies on. That second role raises a question monitoring
alone does not answer: what happens to the evidence when the machine
producing it is compromised, since an attacker who can edit or delete the
logs on a host erases their own trace. This chapter covers collecting logs
so that they are useful and protecting them so that they can be trusted,
which is the integrity property applied to the record itself.

\section{Log sources}

Session~9 named the telemetry that matters, and the same sources are now
captured deliberately (CSA, 2024, p.~128). \emph{Management plane logs}
record control-plane actions, that is, who accessed the infrastructure,
what they changed and when, so an attacker on the console is seen here
first. \emph{Service and application logs} record the activity of
individual services: authentication attempts, API calls, network traffic,
errors. On the red thread these are Nextcloud's own access and audit logs,
the Nginx access log of Session~3, the system journal (\texttt{journalctl})
of Session~2 and the firewall's record of what it dropped. Because an event
that was never logged cannot be investigated, the first rule of logging is
to know what each source captures and to make sure that the
security-relevant sources are actually recorded.

\section{Logs, events and the content of a log line}

The Guidance separates logs from events, and the separation decides how
each is used (CSA, 2024, p.~125). A \emph{log} is a relatively complete
record of activity (create, read, update, delete), detailed and stored
persistently, but often delivered in batches and therefore delayed. An
\emph{event} typically records only a change (create, update, delete), is
not retained unless explicitly saved, and becomes available within seconds,
which is why it can trigger a real-time alert from a service such as AWS
GuardDuty, Microsoft Sentinel or GCP Security Command Center. Events thus
give timely notice that something happened, and logs give the depth to
investigate what it was.

Three lines from the running example, slightly shortened, show what a log
line contains:

\begin{center}
\fcolorbox{kgrey}{klight}{\parbox{0.93\linewidth}{\footnotesize\ttfamily
\# Nginx access log --- a normal request, then a probe\\
192.0.2.44 - - [06/Jul/2026:14:22:07] "GET /apps/files" 200 5120\\
203.0.113.9 - - [06/Jul/2026:14:22:41] "GET /login?user=admin'--" 403\\[3pt]
\# System journal --- a failed SSH login\\
Jul 06 14:23:19 nc-vm sshd[2011]: Failed password for invalid\\
\phantom{xx}user admin from 203.0.113.9 port 51234 ssh2\\[3pt]
\# Azure Activity Log --- a management-plane change\\
operation: Microsoft.Network/networkSecurityGroups/write\\
caller: attacker@example.com \ \ result: Succeeded
}}
\end{center}

Each line carries the same essentials: who (source address or caller), what
(the request or operation), when (the timestamp) and the outcome (the
\texttt{200}, \texttt{403} or \texttt{Succeeded}). The first block shows an
ordinary file fetch followed by an SQL-injection probe against the login,
refused with a \texttt{403}. The second shows a failed SSH login; one is
noise, but fifty a minute is a brute-force attempt. The third is a change to
a network security group on the management plane, one of the key indicators
of Session~9. Each line on its own is mundane; gathered and correlated, such
lines become the timeline of an incident.

In the cloud these sources have concrete names. On Azure the management
plane log is the \emph{Activity Log}, which contains subscription-level
events such as a resource being modified or a VM started; it is enabled by
default, retained for ninety days, and is where forensics after a breach
begins. Per-resource \emph{Resource Logs} (formerly called diagnostic logs)
record the internal activity of a service but are not collected by default,
because a diagnostic setting must first be enabled or an agent installed,
and this gap commonly leaves exactly the logs one wants missing when they
are needed. Both kinds can be routed to \emph{Log Analytics}, an
aggregation-and-query service (queried in KQL) that plays the role of the
SIEM below, or exported to a third-party SIEM through an Event Hub.

\section{Central collection}

Logs scattered across many hosts are hard to search and easy to lose, and
in the cloud, where instances are ephemeral (Session~5), a log left on a
host vanishes with the host. For exactly this reason the Guidance advises
that a workload's logs leave it for a central store as soon as they are
written (CSA, 2024, p.~181).

\begin{definition}[Log aggregation and SIEM]\label{def:logagg}
\textbf{Log aggregation} describes the practice of forwarding logs from many
sources to a central repository for storage, search and analysis. A
\textbf{security information and event management} (SIEM) system describes
such a repository specialised for security: it collects logs across
environments, correlates them to identify potential incidents, and supports
the detection and investigation of Session~9.
\end{definition}

The Guidance calls the resulting architecture a \emph{cascading log
architecture}, in which logs flow from one layer to the next (CSA, 2024,
p.~134). Each environment (development, test, production) sends its logs to
one central store. That store collects everything and passes the subset that
matters for security on to a separate, well-protected audit environment,
where it is kept for the long term, and from there the logs feed a SIEM that
analyses across all environments.

\begin{figure}[htbp]
\centering
\begin{tikzpicture}[font=\footnotesize, node distance=5mm]
  \node[box] (dev)  {Dev};
  \node[box, below=3mm of dev] (test) {Test};
  \node[box, below=3mm of test] (prod) {Prod};
  \node[box, right=16mm of test] (central) {central log\\repository};
  \node[box, right=14mm of central] (audit) {security /\\audit store};
  \node[redbox, right=14mm of audit] (siem) {SIEM};
  \draw[flow] (dev.east)  -- (central.west);
  \draw[flow] (test.east) -- (central.west);
  \draw[flow] (prod.east) -- (central.west);
  \draw[flow] (central) -- (audit) node[midway, above, font=\scriptsize\itshape]{filter};
  \draw[flow] (audit) -- (siem);
\end{tikzpicture}
\caption{Cascading log architecture. Each environment forwards to a central
repository; the security-relevant subset is filtered into a retained audit
store and fed to the SIEM, the same cascade-and-filter as in Session~9, now
applied to storage and cost.}
\label{fig:cascade}
\end{figure}

Two design points from Session~9 carry over. Cascade-and-filter applies
here as cost control (CSA, 2024, pp.~133--135): pertinent alerts and
selected logs go to the security team, bulk raw logs stay cheaply in place,
and rarely-used logs move to lower-cost storage for retention rather than
being discarded. The fast-path/slow-path split applies too, since some logs
feed near-real-time alerting and most feed slower analysis.
\emph{Retention}, that is, how long logs are kept, balances operational
needs, regulatory requirements and cost (CSA, 2024, p.~133): an intrusion
may be discovered months after it began and compliance obligations
(Session~12) set minimum periods, whereas prolonged retention raises storage
cost and has privacy implications. Some log data should also leave the cloud
altogether, both so that it survives a failure of the platform and because
regulators may require it (CSA, 2024, p.~133), and, as the next section
shows, this is an integrity control as well as a resilience measure.

\section{Log integrity}

Everything above assumes that the logs are an honest record, and an
attacker's assumption is the opposite.

\begin{definition}[Log integrity]\label{def:logintegrity}
\textbf{Log integrity} describes the assurance that log records have not been
altered or deleted since they were written. It is the integrity property of
the CIA triad applied to the audit record itself: without it, logs can be
edited to hide an attacker's actions and cannot be trusted as evidence.
\end{definition}

On many intrusions, one of the attacker's first moves after gaining control
of a host is to edit or wipe its logs, removing the record of how they
entered and what they did. If the logs exist only on the compromised
machine, the attacker rewrites the record freely, so the record must be
placed where the compromised host cannot reach it. Three measures follow.
Logs are shipped off-host to the central store of
Definition~\ref{def:logagg} as they are produced, so that a copy exists
beyond the attacker's control the moment the event happens. They are stored
append-only or write-once, so that even with access they can be added to
but not edited or deleted. And access to the log store is restricted under
the least-privilege rules of Session~7. The principle is the one that makes
off-site backups matter in Session~11: a safeguard under the attacker's
control is no safeguard. Centralised collection is therefore itself an
integrity control, because logs that left the host at once are logs the
host's attacker never reached.

\section{Logs as evidence}

The reason integrity is worth this effort is what logs become during an
incident. The Detection and Analysis phase of Session~12 builds a timeline of
the attack from the logs (CSA, 2024, p.~255), and a timeline assembled from
records that might have been tampered with proves nothing. Incident response
therefore speaks of \emph{chain of custody}: evidence is documented and
preserved so that its integrity can be demonstrated later, to management, to
regulators or in legal proceedings (CSA, 2024, pp.~255, 266). A snapshot
(Session~2) and an off-host, append-only log have the same forensic value, a
trustworthy record of a moment the attacker could not revise. On a real
system the logs are therefore worth exactly as much as the confidence that
they still state what happened.

\begin{example}[Protecting the record]\label{ex:logship-trc}
An attacker gains a foothold on the Nextcloud VM. \emph{Threat}: to cover
their tracks, they delete the Nginx access log and edit the system journal,
erasing the evidence of the intrusion. \emph{Risk}: if the only copy is on the
host, the likelihood of successful tampering is high and the impact severe,
because the incident may never be reconstructed and cannot be proven if it
is. \emph{Control}: forward all security-relevant logs off-host to a central
store as they are written, held append-only, with tightly restricted access.
The attacker now controls the host but not the record: the copy that left in
real time survives, the timeline can be rebuilt, and the chain of custody
holds. The control protects the integrity of the evidence rather than its
secrecy, and it is what makes the response of Session~12 possible.
\end{example}

\paragraph{Reading for this chapter.} CSA Security Guidance, Domain~6,
\S6.1.1 (Logs \& Events, pp.~125--126), \S6.2 (Cloud Telemetry Sources,
pp.~127--133), \S6.3.1 (Log Storage \& Retention, p.~133) and \S6.3.2
(Cascading Log Architecture, pp.~134--135); Domain~11, \S11.1.1 (the
incident-response checklist: building the attack timeline and chain of
custody, pp.~254--256) and \S11.3.4 (Cloud System Forensics, pp.~265--267).
Optional deepening: J\o sang, Chapter~16 (Cyber Operations, pp.~337--354), on logging for
detection and forensics. The Microsoft Learn module is \emph{Describe Azure
monitoring} (Azure Monitor, Activity Logs, Log Analytics and diagnostic
settings).

\section{Summary}

Logs are both an operational tool and security evidence, which raises the
question whether the record can be trusted when the host producing it is
compromised. The sources are those of Session~9, management plane logs and
service and application logs, and concretely the records of Nextcloud,
Nginx, the journal and the firewall. Because scattered logs are easily lost
and ephemeral hosts take their logs with them, logs are aggregated centrally
(Definition~\ref{def:logagg}) into a SIEM, with cascade-and-filter for cost,
fast and slow paths for timeliness, and retention set by investigation needs
and compliance. Log integrity (Definition~\ref{def:logintegrity}) requires
that the record leave the host as it is written and be held append-only with
restricted access, so centralised collection is itself an integrity control.
Under chain of custody the logs then become the attack timeline and the
evidence for the response of Session~12 (Example~\ref{ex:logship-trc}).
Session~11 turns from proving harm to surviving it: resilience, RAID and
backup.

\section*{Self-check}
\addcontentsline{toc}{section}{Self-check}

Sketch answers follow the exercises.

\begin{exercise}[The two roles of logs]
Give the two roles logs play, operational and evidential, and one Nextcloud
log source relevant to each.
\end{exercise}

\begin{exercise}[Central collection for ephemeral hosts]
Why is centralised log collection especially important for ephemeral cloud
workloads (Session~5)?
\end{exercise}

\begin{exercise}[Log integrity and the CIA triad]
Define log integrity in terms of the CIA triad, and explain why logs kept
only on a host fail it once the host is compromised.
\end{exercise}

\begin{exercise}[Protecting log integrity]
Name three concrete measures that protect log integrity, and state why
centralised collection is an integrity control and not only a convenience.
\end{exercise}

\begin{exercise}[Chain of custody]
What is chain of custody, and why does it make off-host, append-only logging
matter for the incident response of Session~12?
\end{exercise}

\paragraph{Sketch answers.}
{\small
(1)~Operational: understanding the system and reconstructing behaviour, for
example the Nginx access log or the system journal. Evidential: the timeline
and proof of who did what during an incident, for example Nextcloud's audit
log and the management plane logs. (2)~Ephemeral instances disappear with
demand, taking locally stored logs with them; forwarding logs to a central
store as they are written preserves them beyond the life of the host.
(3)~Log integrity is the integrity property of CIA applied to the audit
record, the assurance that it has not been altered or deleted. Logs kept only
on the host fail it because an attacker who controls the host can edit or
wipe them, destroying the record of their own actions. (4)~Ship logs off-host
as they are produced; store them append-only or write-once; restrict who may
access the log store (least privilege). Central collection is an integrity
control because logs that left the host immediately are beyond the reach of
the host's attacker. (5)~Chain of custody is the documented preservation of
evidence so that its integrity can be demonstrated later (to management,
regulators or a court); off-host, append-only logs provide a trustworthy,
untampered record, which is what lets the attack timeline in Session~12's
Detection and Analysis phase stand as proof.}


